\chapter{Haemorrhoidectomy}
\index{Haemorrhoidectomy}

\section*{Definition and Overview}
Haemorrhoidectomy is a surgical procedure to remove haemorrhoids (piles), the swollen veins in the lower rectum and anus. It is recommended for haemorrhoids that are large, prolapse (protrude) through the anus, bleed repeatedly, or do not respond to conservative treatments and less invasive procedures. The operation removes the haemorrhoidal tissue, usually under general or spinal anaesthetic, and is performed as day surgery or with a short hospital stay. Haemorrhoidectomy is highly effective, with the lowest recurrence rate of all haemorrhoid treatments, but it has a more painful recovery than other options. Most people return to normal activities within two to four weeks.

\section*{Epidemiology}
Haemorrhoids affect around 1 in 20 people at any time, and up to half of adults have them at some stage. Only a minority need surgery: most are managed with diet, fibre, creams, and office procedures such as rubber band ligation. Haemorrhoidectomy is performed for the most troublesome cases, most often for grade 3 and 4 haemorrhoids, which prolapse and do not retract on their own. It is more common with increasing age, and the operation is equally performed in men and women, including during pregnancy planning.

\section*{Aetiology and Pathophysiology}
Haemorrhoids are cushions of vascular tissue in the anal canal that can enlarge and become symptomatic.

\begin{itemize}
    \item \textbf{Normal anatomy:} everyone has anal cushions that help maintain continence
    \item \textbf{Enlargement:} straining, constipation, pregnancy, obesity, and prolonged sitting increase pressure and cause the cushions to enlarge and prolapse
    \item \textbf{Grading:} grade 1 haemorrhoids bleed but do not prolapse; grades 2--4 prolapse increasingly and eventually cannot be pushed back in
    \item \textbf{Symptoms:} bleeding, discomfort, itching, and prolapse
    \item \textbf{Surgery rationale:} removing the haemorrhoidal tissue relieves symptoms and prevents recurrence
    \item \textbf{Stapled haemorrhoidectomy:} an alternative technique that repositions rather than removes tissue
\end{itemize}

\section*{Clinical Features}
Surgery is considered for people with specific symptoms.

\begin{itemize}
    \item Haemorrhoids that prolapse and require manual pushing back, or that cannot be replaced
    \item Repeated bleeding from haemorrhoids
    \item Pain and discomfort from thrombosed or strangulated haemorrhoids
    \item Symptoms that persist despite diet, fibre, and less invasive treatment
    \item Haemorrhoids that interfere with daily life, hygiene, or work
    \item Failure or unsuitability of rubber band ligation and other office procedures
    \item Large external haemorrhoids causing hygiene problems
\end{itemize}

\section*{Differential Diagnosis}
Before surgery, other causes of anal symptoms must be excluded.

\begin{itemize}
    \item Anal fissures, which cause pain and bleeding on defecation
    \item Anal fistula and abscess
    \item Rectal polyps and rectal cancer, which can bleed
    \item Inflammatory bowel disease, including Crohn disease
    \item Rectal prolapse
    \item Itchy skin conditions of the anus
    \item Other causes of rectal bleeding, which require investigation
\end{itemize}

\section*{When to Seek Medical Care}
People with persistent or severe anal symptoms should see their GP, who can examine, arrange any necessary tests, and refer to a colorectal or general surgeon. Rectal bleeding should always be investigated, especially in people over 50.

\textbf{Contact your local emergency services for:}
\begin{itemize}
    \item Heavy rectal bleeding or bleeding that will not stop
    \item Sudden severe anal pain with a tender, swollen lump (thrombosed or strangulated haemorrhoid)
    \item Fever, spreading redness, or severe pain after surgery, which may indicate infection
    \item Inability to pass urine after surgery
    \item Signs of a serious complication, including heavy bleeding or difficulty breathing
\end{itemize}

\section*{Diagnosis and Investigations}
Pre-operative assessment confirms the diagnosis and fitness for surgery.

\begin{itemize}
    \item \textbf{Clinical examination:} inspection and digital rectal examination
    \item \textbf{Proctoscopy:} visual inspection of the anal canal and lower rectum
    \item \textbf{Colonoscopy:} recommended when there is bleeding and risk factors for bowel cancer, to exclude other causes
    \item \textbf{Full blood count:} to assess anaemia from chronic bleeding
    \item \textbf{General health assessment:} fitness for anaesthesia
    \item \textbf{Discussion of options:} the surgeon explains the procedure, expected recovery, and alternatives
\end{itemize}

\section*{Management}
Haemorrhoidectomy is one of several options, chosen when surgery is most appropriate.

\begin{itemize}
    \item \textbf{Open or closed haemorrhoidectomy:} the standard operation, removing the haemorrhoidal tissue
    \item \textbf{Stapled haemorrhoidectomy:} a less painful alternative for internal haemorrhoids, using a stapling device to reposition tissue
    \item \textbf{Laser and other techniques:} used in some centres
    \item \textbf{Anaesthesia:} general or spinal anaesthetic
    \item \textbf{Recovery:} pain is expected for one to two weeks, managed with painkillers, and most people take one to two weeks off work
    \item \textbf{Bowel care:} stool softeners and fibre prevent constipation and straining after surgery
    \item \textbf{Follow-up:} review to check healing and discuss results
    \item \textbf{Lifestyle measures:} fibre, fluids, and avoiding straining reduce recurrence
\end{itemize}

\section*{Complications}
\begin{itemize}
    \item Pain, which is the most common complaint after surgery
    \item Bleeding, including reactionary bleeding in the first 24 hours and secondary bleeding at one to two weeks
    \item Infection and abscess formation
    \item Urinary retention, requiring catheterisation in some cases
    \item Anal stenosis, narrowing of the anal canal from scarring
    \item Incontinence or urgency, which is usually temporary
    \item Recurrence of haemorrhoids, which occurs in a small proportion
    \item Complications of anaesthesia
\end{itemize}

\section*{Prognosis and Outcomes}
Haemorrhoidectomy is the most effective treatment for haemorrhoids, with long-term cure in over 90% of people and the lowest recurrence rate of any treatment. Recovery is uncomfortable but predictable, and most people are back to normal within two to four weeks. Complications are uncommon with modern technique and good aftercare. The vast majority of people who have the operation are satisfied with the result and are free of the bleeding, prolapse, and discomfort that prompted surgery.

\section*{Prevention and Self-Care}
\begin{itemize}
    \item Prevent constipation with a high-fibre diet, plenty of fluids, and regular exercise
    \item Avoid straining on the toilet, and do not sit for prolonged periods
    \item Go to the toilet when you feel the urge, rather than delaying
    \item After surgery, take stool softeners and pain relief as directed
    \item Keep the anal area clean and dry, and use warm baths for comfort
    \item Avoid heavy lifting and strenuous activity for the advised period
    \item Attend follow-up appointments
    \item Report heavy bleeding, fever, or severe pain promptly
    \item Continue fibre and fluid intake long term to prevent recurrence
\end{itemize}

\section*{Sources and Further Reading}
The following reputable sources provide further information for healthcare students and patients seeking reliable, current advice about haemorrhoidectomy.

\begin{itemize}
    \item Royal Australasian College of Surgeons. \textit{Haemorrhoidectomy patient information.} https://www.surgeons.org/
    \item Australian Society of Colorectal Surgery. \textit{Haemorrhoids and their treatment.} https://www.colorectalsurgeons.com.au/
    \item Healthdirect Australia. \textit{Haemorrhoids (piles) treatment.} https://www.healthdirect.gov.au/
    \item Gut Foundation of Australia. \textit{Haemorrhoids information.} https://gutfoundation.com.au/
    \item NHS. \textit{Haemorrhoidectomy.} https://www.nhs.uk/
    \item American Society of Colon and Rectal Surgeons. \textit{Hemorrhoid surgery.} https://fascrs.org/
\end{itemize}
